% ======================================
% Appendix C: Source Code
% ======================================

This appendix presents the key source code files that implement the core functionality of the Quantum Virtual Machine (QVM) system. The complete source code repository is available at:

\begin{center}
\textbf{\url{https://github.com/qayum/quantum-virtual-machine}}
\end{center}

\vspace{0.5cm}

\textbf{Project Information:}
\begin{itemize}
\item \textbf{Student 1:} Muhammad Abdul Qayyoum (CIIT/SP23-BCS-033/VHR)
\item \textbf{Student 2:} Ahmad Faraz (CIIT/SP23-BCS-007/VHR)
\item \textbf{Supervisor:} Assistant Prof. Mr. Farhad Mubashir
\item \textbf{Repository:} \url{https://github.com/qayum/quantum-virtual-machine}
\item \textbf{Branch:} main
\item \textbf{Language:} Python 3.10+
\item \textbf{Dependencies:} NumPy, Matplotlib, Lark, FastAPI
\end{itemize}

\vspace{0.5cm}

The following sections present eight key source files that demonstrate the system's architecture and implementation. These files represent the core modules: Intermediate Representation, Parsing, Simulation, Transpilation, User Interface, Hardware Architecture, and Gate Decomposition.

\clearpage

\subsection{File 1: Intermediate Representation (ir.py)}

\textbf{Purpose:} Defines the QuantumCircuit class that serves as the unified internal representation for quantum circuits across all input formats.

\textbf{Key Features:}
\begin{itemize}
\item QuantumCircuit class with qubit and classical register management
\item Support for conditional operations based on classical bit values
\item Control flow primitives (labels and jumps) for loops and branching
\item Classical bit operations (AND, OR, XOR, NOT, assignment)
\item Extensible operation dictionary structure for gate metadata
\end{itemize}

\textbf{Location:} \texttt{src/qvm/ir.py}

\begin{lstlisting}[language=Python, caption={Intermediate Representation Module}]
# src/qvm/ir.py

"""
Intermediate Representation (IR) for Quantum Circuits.
Extends the IR to support OpenQASM 3.0 classical registers and conditional operations.
"""

from typing import List, Dict, Optional, Any

class QuantumCircuit:
    def __init__(self, num_qubits: int):
        if not isinstance(num_qubits, int) or num_qubits <= 0:
            raise ValueError("Number of qubits must be a positive integer.")
        self.num_qubits = num_qubits
        self.operations = []  # List of dictionaries: gate, qubits, params, condition, target_bit
        self.classical_registers: Dict[str, int] = {}  # name -> size

    def add_classical_register(self, name: str, size: int):
        """Declares a classical bit register."""
        if name in self.classical_registers:
            raise ValueError(f"Classical register '{name}' already exists.")
        self.classical_registers[name] = size

    def add_operation(self, gate_name: str, qubits: list, params: list = None, 
                     condition: dict = None, target_bit: tuple = None, 
                     duration: str = None, label: str = None, jump_to: str = None, 
                     classical_op: dict = None):
        """
        Adds a quantum or classical operation to the circuit.

        Args:
            gate_name (str): The name (e.g., "h", "measure", "classical_op").
            qubits (list): Target quantum bits.
            params (list, optional): Gate parameters.
            condition (dict, optional): {"register": str, "index": int, "value": int}
            target_bit (tuple, optional): (register_name, index) for results.
            duration (str, optional): Timing string.
            label (str, optional): Label for jumps.
            jump_to (str, optional): Target label for jumps.
            classical_op (dict, optional): {"op": str, "target": tuple, "args": list}
        """
        if condition:
            reg = condition.get("register")
            if reg not in self.classical_registers:
                raise ValueError(f"Unknown classical register in condition: {reg}")
            if not (0 <= condition.get("index", 0) < self.classical_registers[reg]):
                raise ValueError(f"Index out of bounds for classical register '{reg}'")

        operation = {
            "name": gate_name,
            "qubits": qubits if qubits is not None else [],
            "params": params if params is not None else [],
            "condition": condition,
            "target_bit": target_bit,
            "duration": duration,
            "label": label,
            "jump_to": jump_to,
            "classical_op": classical_op
        }
        self.operations.append(operation)

    def __str__(self):
        s = f"QuantumCircuit(num_qubits={self.num_qubits}, registers={self.classical_registers})\n"
        for op in self.operations:
            if op["name"] == "label":
                s += f"  LABEL {op['label']}:\n"
                continue
            if op["name"] == "jump":
                cond_str = f" IF {op['condition']}" if op['condition'] else ""
                s += f"  JUMP {op['jump_to']}{cond_str}\n"
                continue
            if op["name"] == "classical_op":
                s += f"  CLASSICAL {op['classical_op']['target']} = {op['classical_op']['op']} {op['classical_op']['args']}\n"
                continue
            
            cond_str = f" IF {op['condition']}" if op['condition'] else ""
            target_str = f" -> {op['target_bit']}" if op['target_bit'] else ""
            dur_str = f" [{op['duration']}]" if op['duration'] else ""
            s += f"  {op['name']}{dur_str} {op['qubits']}{cond_str}{target_str}\n"
        return s
\end{lstlisting}

\clearpage

\subsection{File 2: Hardware Architecture (architecture.py)}

\textbf{Purpose:} Defines data structures for representing target quantum hardware topologies and connectivity constraints.

\textbf{Key Features:}
\begin{itemize}
\item TargetArchitecture class for hardware specification
\item Bidirectional connectivity validation
\item Pre-defined architectures (linear chain, fully connected)
\item Native gate set specification
\item Connectivity checking for qubit pairs
\end{itemize}

\textbf{Location:} \texttt{src/qvm/architecture.py}

\begin{lstlisting}[language=Python, caption={Hardware Architecture Module}]
# src/qvm/architecture.py

"""
Defines data structures for representing target quantum hardware architectures.
"""

from typing import List, Set, Tuple

class TargetArchitecture:
    """
    Represents the constraints of a target quantum hardware, including
    qubit connectivity and the native gate set.
    """
    def __init__(self, name: str, num_qubits: int, connectivity: Set[Tuple[int, int]], 
                 native_gates: Set[str]):
        """
        Initializes a new TargetArchitecture.

        Args:
            name (str): The name of the architecture (e.g., "Linear-4", "IBM-Q-5").
            num_qubits (int): The total number of qubits in the architecture.
            connectivity (Set[Tuple[int, int]]): A set of tuples where each tuple represents
                                                  a physical connection between two qubits.
                                                  Connections are assumed to be bidirectional.
            native_gates (Set[str]): A set of gate names that are natively supported by the hardware.
        """
        if not isinstance(num_qubits, int) or num_qubits <= 0:
            raise ValueError("Number of qubits must be a positive integer.")
        
        self.name = name
        self.num_qubits = num_qubits
        self.connectivity = self._validate_and_normalize_connectivity(connectivity, num_qubits)
        self.native_gates = native_gates

    def _validate_and_normalize_connectivity(self, connectivity: Set[Tuple[int, int]], 
                                            num_qubits: int) -> Set[Tuple[int, int]]:
        """Validates and ensures bidirectionality of connectivity."""
        normalized = set()
        for edge in connectivity:
            if not (isinstance(edge, tuple) and len(edge) == 2 and
                    isinstance(edge[0], int) and isinstance(edge[1], int) and
                    0 <= edge[0] < num_qubits and 0 <= edge[1] < num_qubits):
                raise ValueError(f"Invalid edge in connectivity: {edge}")
            # Add both (u, v) and (v, u) to ensure bidirectionality and easy lookup
            normalized.add(tuple(sorted(edge)))
        return normalized

    def is_connected(self, qubit1: int, qubit2: int) -> bool:
        """Checks if two physical qubits are directly connected."""
        return tuple(sorted((qubit1, qubit2))) in self.connectivity

    def __str__(self):
        return f"TargetArchitecture(name='{self.name}', num_qubits={self.num_qubits}, native_gates={self.native_gates})"

# Pre-defined architectures for convenience
def get_linear_architecture(num_qubits: int) -> TargetArchitecture:
    """Creates a linear chain architecture."""
    if num_qubits < 2:
        connectivity = set()
    else:
        connectivity = {(i, i + 1) for i in range(num_qubits - 1)}
    
    # A common, minimal native gate set
    native_gates = {"id", "rz", "sx", "x", "cx"}

    return TargetArchitecture(f"Linear-{num_qubits}", num_qubits, connectivity, native_gates)

def get_fully_connected_architecture(num_qubits: int) -> TargetArchitecture:
    """Creates a fully connected (all-to-all) architecture."""
    connectivity = set()
    if num_qubits >= 2:
        for i in range(num_qubits):
            for j in range(i + 1, num_qubits):
                connectivity.add((i, j))

    native_gates = {"id", "rz", "sx", "x", "cx"}

    return TargetArchitecture(f"FullyConnected-{num_qubits}", num_qubits, connectivity, native_gates)
\end{lstlisting}

\clearpage

\subsection{File 3: Gate Decomposer (decomposer.py)}

\textbf{Purpose:} Decomposes complex quantum gates into sequences of simpler native gates supported by the target hardware.

\textbf{Key Features:}
\begin{itemize}
\item Toffoli (CCX) gate decomposition into single and two-qubit gates
\item Preservation of conditional execution metadata
\item Support for OpenQASM 3.0 control flow primitives
\item Extensible decomposition rule system
\item Circuit-level decomposition with classical register preservation
\end{itemize}

\textbf{Location:} \texttt{src/qvm/decomposer.py}

\begin{lstlisting}[language=Python, caption={Gate Decomposition Module}]
# src/qvm/decomposer.py

"""
Decomposes complex quantum gates into a sequence of simpler, native gates.
Updated to preserve OpenQASM 3.0 classical metadata and control flow.
"""

from src.qvm.ir import QuantumCircuit

class Decomposer:
    """
    A simple decomposer that can break down specific complex gates.
    """
    def __init__(self, native_gates: set):
        self.native_gates = native_gates

    def decompose_operation(self, op: dict) -> list:
        """
        Decomposes a single gate operation if it's not in the native gate set.
        Returns a list of simpler operations.
        """
        gate_name = op["name"]
        
        # If it's a non-gate op (classical, label, etc.), it's considered native
        if gate_name in ["classical_op", "label", "jump", "delay", "measure"]:
            return [op]

        if gate_name in self.native_gates:
            return [op]

        if gate_name == "toffoli" or gate_name == "ccx":
            return self._decompose_toffoli(op)
        
        raise ValueError(f"No decomposition rule available for gate: {gate_name}")

    def _decompose_toffoli(self, op: dict) -> list:
        qubits = op["qubits"]
        if len(qubits) != 3:
            raise ValueError("Toffoli gate must act on 3 qubits.")
        c1, c2, t = qubits[0], qubits[1], qubits[2]
        
        import numpy as np
        pi_4 = np.pi / 4

        # Preserve condition if the CCX was conditional
        cond = op.get("condition")

        decomposition = [
            {"name": "h", "qubits": [t], "params": [], "condition": cond},
            {"name": "cx", "qubits": [c2, t], "params": [], "condition": cond},
            {"name": "rz", "qubits": [t], "params": [-pi_4], "condition": cond},
            {"name": "cx", "qubits": [c1, t], "params": [], "condition": cond},
            {"name": "rz", "qubits": [t], "params": [pi_4], "condition": cond},
            {"name": "cx", "qubits": [c2, t], "params": [], "condition": cond},
            {"name": "rz", "qubits": [t], "params": [-pi_4], "condition": cond},
            {"name": "cx", "qubits": [c1, t], "params": [], "condition": cond},
            {"name": "rz", "qubits": [c2], "params": [pi_4], "condition": cond},
            {"name": "rz", "qubits": [t], "params": [pi_4], "condition": cond},
            {"name": "h", "qubits": [t], "params": [], "condition": cond},
            {"name": "cx", "qubits": [c1, c2], "params": [], "condition": cond},
            {"name": "rz", "qubits": [c1], "params": [pi_4], "condition": cond},
            {"name": "rz", "qubits": [c2], "params": [-pi_4], "condition": cond},
            {"name": "cx", "qubits": [c1, c2], "params": [], "condition": cond},
        ]
        return decomposition

    def decompose_circuit(self, circuit: QuantumCircuit) -> QuantumCircuit:
        """
        Decomposes all non-native gates in a circuit.
        """
        decomposed_circuit = QuantumCircuit(circuit.num_qubits)
        decomposed_circuit.classical_registers = circuit.classical_registers.copy()
        
        for op in circuit.operations:
            decomposed_ops = self.decompose_operation(op)
            for new_op in decomposed_ops:
                decomposed_circuit.add_operation(
                    new_op["name"], 
                    new_op.get("qubits", []), 
                    params=new_op.get("params", []),
                    condition=new_op.get("condition"),
                    target_bit=new_op.get("target_bit"),
                    duration=new_op.get("duration"),
                    label=new_op.get("label"),
                    jump_to=new_op.get("jump_to"),
                    classical_op=new_op.get("classical_op")
                )
        return decomposed_circuit
\end{lstlisting}

\clearpage

\textbf{Note:} Due to space constraints, the remaining source files (Parser, Simulator, Transpiler, CLI, and QASM3 Parser) are available in the complete repository at the URL provided above. These files implement:

\begin{itemize}
\item \textbf{parser.py} (150 lines): OpenQASM 2.0 parsing with syntax validation
\item \textbf{simulator.py} (200 lines): Statevector simulation with classical memory and control flow
\item \textbf{transpiler.py} (200 lines): Hardware-aware compilation with Greedy BFS and SABRE routing
\item \textbf{cli.py} (100 lines): Command-line interface with argument parsing and visualization
\item \textbf{qasm3\_parser.py} (300 lines): OpenQASM 3.0 parser using Lark grammar with control flow support
\end{itemize}

\vspace{1cm}

\textbf{Repository Structure:}
\begin{verbatim}
quantum-virtual-machine/
├── src/
│   ├── qvm/
│   │   ├── ir.py                 # Intermediate Representation
│   │   ├── parser.py             # OpenQASM 2.0 Parser
│   │   ├── qasm3_parser.py       # OpenQASM 3.0 Parser
│   │   ├── simulator.py          # Statevector Simulator
│   │   ├── transpiler.py         # Hardware-Aware Transpiler
│   │   ├── decomposer.py         # Gate Decomposer
│   │   ├── architecture.py       # Hardware Topology
│   │   ├── visual.py             # Visualization
│   │   ├── cli.py                # Command-Line Interface
│   │   └── server.py             # Web API Server
│   └── examples/
│       └── full_pipeline.py      # Example Usage
├── tests/                        # Unit Tests
├── docs/                         # Documentation
├── requirements.txt              # Python Dependencies
└── README.md                     # Setup Instructions
\end{verbatim}

\vspace{0.5cm}

\textbf{Setup Instructions:}
\begin{enumerate}
\item Clone repository: \texttt{git clone https://github.com/qayum/quantum-virtual-machine.git}
\item Install dependencies: \texttt{pip install -r requirements.txt}
\item Run CLI: \texttt{python -m src.qvm.cli examples/bell\_state.qasm}
\item Run Web Server: \texttt{python -m src.qvm.server}
\item Run Tests: \texttt{pytest tests/}
\end{enumerate}

\vspace{0.5cm}

\textbf{System Requirements:}
\begin{itemize}
\item Python 3.10 or higher
\item NumPy 1.21+ (linear algebra operations)
\item Matplotlib 3.5+ (visualization)
\item Lark 1.1+ (OpenQASM 3.0 parsing)
\item FastAPI 0.95+ (web API)
\item Uvicorn (ASGI server)
\end{itemize}
